% DRAFT: This report is a draft  and needs review if anyone available.
\documentclass[sigconf,nonacm]{acmart}

\settopmatter{printacmref=false}
\renewcommand\footnotetextcopyrightpermission[1]{}
\setcopyright{none}

\begin{document}

\title{Honey Encryption for Structured English Messages}

% TODO: Add  email addresses and Nidya  full name when available, also finish discussion.
\author{Jackson Murphy}
\affiliation{%
  \institution{Applied Cryptography}
}

\author{Kush Patel}
\affiliation{%
  \institution{Applied Cryptography}
}

\author{Mourtalla Diop}
\affiliation{%
  \institution{Applied Cryptography}
}

\author{Nidya}
\affiliation{%
  \institution{Applied Cryptography}
}

\begin{abstract}
Honey encryption aims to make an incorrect decryption key produce a plausible
message rather than an obvious error. This project explores that idea in a
small, controlled setting: a prototype maps 256-bit seeds to one of 112
structured meeting messages and uses a password-derived mask to protect the
seed. We evaluate the uniform message encoder, a separate weighted encoder
using a synthetic prior, and the behavior of incorrect passwords and simple
side information. The results show that the encoders produce messages in the
specified format and approximately follow their configured seed
distributions. They do not establish that the outputs are indistinguishable
from real messages or that the prototype is secure. This report describes the
design, implementation, tests, results, limitations, and possible extensions.
\end{abstract}

\keywords{honey encryption, distribution-transforming encoder, cryptography, evaluation}

\maketitle

\section{Introduction}
\label{sec:introduction}

Honey encryption is a cryptographic approach intended to change the outcome
of password guessing: an incorrect password should yield a plausible-looking
decoy rather than a clear failure. Its central challenge is constructing a
distribution-transforming encoder (DTE) that maps messages to seeds in a way
that reflects the expected message distribution
\cite{juels2014honey}.

This project studies that challenge in a deliberately restricted domain. The
prototype supports 112 messages of the form
\emph{``Meet me at [location] on [day] at [time].''} The project asks whether
a DTE for this finite message set can produce correctly formatted outputs for
arbitrary seeds, how closely sampled seeds follow the encoder's target
distribution, and how known message fields reduce the set of possible
messages.

The project implements a uniform DTE and an additional weighted DTE based on
a hand-authored synthetic prior. It also measures the output of the
password-wrapped prototype under incorrect guesses and includes a separate
authenticated-encryption baseline for contrast. Unlike work targeting
general-language encryption or a deployed security construction, this is a
small educational prototype with explicit assumptions and measured
limitations. It does not claim formal honey-encryption security.

\section{Background}
\label{sec:background}

\subsection{Honey encryption and distribution-transforming encoders}

In honey encryption, a DTE encodes a message as a seed, and a password-based
cryptographic layer protects that seed. Decryption with a wrong password
produces a different seed, which the DTE decodes as another message. The
intended benefit depends on decoy messages being plausible under a suitable
message distribution; merely producing a string in the expected format is
not enough \cite{juels2014honey}.

This prototype uses a finite message space rather than unrestricted English.
Its uniform DTE assigns equal target probability to each of the 112 supported
sentences. The weighted extension uses synthetic integer weights over
locations, weekdays, and times. Those weights are illustrative and were not
learned from real communications.

\subsection{Project resources}

The project used the original honey-encryption publication as background
\cite{juels2014honey}, the ACM Conference Proceedings Primary Article
template for this report, and its own Python prototype, unit tests, and
captured experiment outputs. The AES-GCM baseline is included to contrast
authenticated encryption behavior; the GCM mode is specified by NIST
\cite{nist2007gcm}. No real message corpus or human-password dataset was used.

\section{System Design and Methodology}
\label{sec:design}

\subsection{Prototype architecture}

The prototype has two main stages. First, the DTE encodes a supported
plaintext sentence into a 256-bit seed. Second, the password layer derives a
mask from the password and a random salt using PBKDF2-HMAC-SHA256 with
200{,}000 iterations, then XORs the mask with the seed. Decryption derives
the mask again, recovers a seed, and passes it to the DTE decoder. Because
the decoder is defined over the seed space, a valid-length wrong password
also decodes to one of the supported sentences.

The weighted DTE is currently a standalone component and is not connected to
the password wrapper. The AES-256-GCM/scrypt implementation is a separate
baseline, not a component of the honey-encryption-inspired construction.
Unlike the prototype's unauthenticated XOR layer, the baseline reports
authentication failure when decryption uses an incorrect password.

\subsection{Evaluation methodology}

The evaluation covers several distinct questions:
\begin{itemize}
  \item Sample 1{,}120{,}000 uniformly random 256-bit seeds and compare the
  uniform DTE output frequencies with the uniform target distribution.
  \item Apply 560 distinct random incorrect-password guesses to one fixed
  ciphertext and salt, then summarize the resulting message frequencies.
  \item Sample 560{,}000 random seeds for the weighted DTE and compare the
  observed frequencies with its declared synthetic target probabilities.
  \item Count the messages compatible with selected known fields, and run a
  small synthetic password-rank pilot.
  \item Run unit tests for DTE boundaries, message round trips, serialization,
  weighted-DTE behavior, and the authenticated-encryption baseline.
\end{itemize}

For distribution experiments, the project reports Pearson chi-square
statistics, p-values, and total variation distance (TVD). These are
descriptive sample comparisons, not security proofs. In particular, the
incorrect-password experiment has an expected count of only five per message,
so the chi-square approximation is fragile.

\section{Implementation and Testing/Results}
\label{sec:results}

\subsection{Implementation}

The uniform prototype defines four locations, seven weekdays, and four times,
for a total of 112 structured sentences. The uniform encoder partitions the
256-bit seed space into intervals corresponding to those messages. The
password layer uses a 16-byte salt, PBKDF2-HMAC-SHA256 with 200{,}000
iterations, and XOR; it does not provide an authentication tag or integrity
protection.

The extension adds an integer-weight DTE with a synthetic product prior,
evaluation scripts, side-information counts, and a small synthetic
dictionary pilot. A separate AES-256-GCM/scrypt baseline provides a
conventional authenticated-encryption comparison. The included test suite
checks the uniform and weighted encoder boundaries, round trips for all
supported messages, probability normalization, JSON behavior, independence
from teaching metadata, and authentication behavior of the baseline.

\subsection{Test results}

In the recorded test run, all 10 unit tests passed. Correct-password
recovery also passed in the recorded incorrect-password experiment. The
AES-GCM baseline recovered a message with the correct password and rejected
an incorrect password through authentication failure. These tests establish
expected behavior for the implementation; they do not establish cryptographic
security.

\subsection{Distribution and side-information results}

Table~\ref{tab:distribution-results} summarizes the recorded distribution
experiments. For random seeds, the uniform DTE produced a TVD of 0.004047
from the uniform target. The weighted DTE produced a TVD of 0.005621 from
its synthetic target. The corresponding p-values indicate that these samples
did not provide evidence against their configured distributions; they are
not proof that either target is realistic.

\begin{table}[t]
  \caption{Recorded seed-distribution experiment results.}
  \label{tab:distribution-results}
  \begin{tabular}{lrrrr}
    \toprule
    Experiment & Trials & $\chi^2$ & $p$-value & TVD \\
    \midrule
    Uniform seeds & 1{,}120{,}000 & 117.444 & 0.319549 & 0.004047 \\
    Weighted seeds & 560{,}000 & 111.004 & 0.482054 & 0.005621 \\
    \bottomrule
  \end{tabular}
\end{table}

For the 560 incorrect-password trials, the recorded chi-square statistic was
100.000 with a p-value of 0.763939 and a TVD of 0.162500. This experiment
used random high-entropy guesses against one fixed ciphertext and salt. It
does not model human-selected passwords or a realistic attacker strategy,
and its small expected count per message makes the chi-square approximation
fragile.

The supported-message space has 112 candidates when no fields are known.
Knowing a location leaves 28 candidates; a day leaves 16; a time leaves 28;
location and day together leave 4; and day and time together also leave 4.
These are exact counts for the prototype's message model, but they do not
identify which candidate is genuine.

The synthetic dictionary pilot used 12 trials and a hand-written list of 25
candidate passwords. It recorded a mean true-password rank of 10.25 and a
median rank of 11. The list is not representative of human passwords, and
the pilot offers no oracle for deciding whether a decrypted message is the
genuine plaintext.

\section{Discussion}
\label{sec:discussion}


\bibliographystyle{ACM-Reference-Format}
\bibliography{references}

\end{document}
